\section{Methodology}

\subsection{System overview and pipeline}
\begin{bluenote}
Present the system overview and the pipeline of the thesis. It should be clear how the digital twin
is used to generate candidate plans and how the Monte Carlo evaluation protocol is used to assess
them.
\end{bluenote}

The project is organized in four different main modules, each one corresponding to one step of the
pipeline :

\begin{table}[htbp]
\centering
\begin{tabularx}{\linewidth}{@{}l>{\raggedright\arraybackslash}X>{\raggedright\arraybackslash}X@{}}
\toprule
Module & Main role & Main outputs \\
\midrule
\texttt{scenario\_generator} & Generates synthetic traffic demand and initial routes & \texttt{routes.xml}, OD matrices \\
\texttt{deviation\_plan\_maker} & Converts the network into a graph and calculates the deviation plan(s) & JSON of candidates deviation plans \\
\texttt{monte\_carlo\_simulation} & SUMO simulation with multiple runs with random seeds & \texttt{run\_\#} folders, \texttt{summary.json} \\
\texttt{evaluator} & Statistical analysis and multi-criteria ranking & \texttt{evaluation.json} \\
\bottomrule
\end{tabularx}
\caption{Functional organization of the pipeline}
\end{table}

\begin{figure}[htbp]
\centering
\begin{tikzpicture}[
  module/.style={
    draw, rounded corners=5pt, fill=blue!7, thick,
    minimum width=5.8cm, minimum height=1.0cm,
    align=center, font=\small\ttfamily
  },
  artifact/.style={
    draw, rounded corners=3pt, fill=gray!8,
    minimum width=4.6cm, minimum height=0.75cm,
    align=center, font=\scriptsize\ttfamily
  },
  inputbox/.style={
    draw, dashed, rounded corners=3pt, fill=orange!7,
    minimum width=5.8cm, minimum height=0.85cm,
    align=center, font=\small
  },
  arr/.style={-{Stealth[length=6pt]}, thick},
  side/.style={-{Stealth[length=5pt]}, thick, gray!70},
  node distance=0.75cm
]
  \node[inputbox] (inp)
    {\texttt{.net.xml} \enspace \texttt{taz.xml} \enspace \texttt{config.toml}};

  \node[module, below=0.85cm of inp]  (sg)  {scenario\_generator};
  \node[module, below=0.75cm of sg]   (dpm) {deviation\_plan\_maker};
  \node[module, below=0.75cm of dpm]  (mcs) {monte\_carlo\_simulation};
  \node[module, below=0.75cm of mcs]  (ev)  {evaluator};

  \node[artifact, right=1.6cm of sg]  (sg_out)  {\texttt{routes.xml}, OD matrices};
  \node[artifact, right=1.6cm of dpm] (dpm_out) {\texttt{plans.json}};
  \node[artifact, right=1.6cm of mcs] (mcs_out) {\texttt{run\_\#/}, \texttt{summary.json}};
  \node[artifact, right=1.6cm of ev]  (ev_out)  {\texttt{evaluation.json}};

  \draw[arr] (inp) -- (sg);
  \draw[arr] (sg)  -- (dpm);
  \draw[arr] (dpm) -- (mcs);
  \draw[arr] (mcs) -- (ev);

  \draw[side] (sg)  -- (sg_out);
  \draw[side] (dpm) -- (dpm_out);
  \draw[side] (mcs) -- (mcs_out);
  \draw[side] (ev)  -- (ev_out);
\end{tikzpicture}
\caption{Four-module pipeline. Inputs (dashed box) flow top-to-bottom through the four modules. Each module's main output artefacts are shown to the right.}
\label{fig:pipeline}
\end{figure}

\subsection{Scenario generation}
The role of the scenario generator is to create demand that is :
\begin{itemize}
\item \textbf{network-aware} because routes are built on the real road network
\item \textbf{zone-aware} because demand is expressed at the TAZ level
\item \textbf{temporally structured} because daily traffic should exhibit morning and evening peaks
\item \textbf{directionally realistic} because peak periods should concentrate trips from
residential areas toward employment centers and back.
\end{itemize}

\subsubsection{Module overview}

\begin{table}[htbp]
  \centering
  \begin{tabularx}{\linewidth}{@{}l>{\raggedright\arraybackslash}X@{}}
  \toprule
  Module & Main role \\
  \midrule
  \texttt{generator.py} & command-line entry point, loads \texttt{TOML} config and launches generation \\
  \texttt{config\_manager.py} & typed access to configuration values and path resolution \\
  \texttt{zones.py} & TAZ loading, edge filtering, zone mass and endpoint sampling \\
  \texttt{od\_model.py} & OD model, temporal profile \\
  \texttt{route\_library.py} & TAZ-to-TAZ impedance estimation and candidate-route generation with caching \\
  \texttt{vehicle\_sampler.py} & conversion from hourly OD cells to individual vehicles with route choice \\
  \texttt{exports.py} & writing \texttt{routes.xml}, \texttt{out\_od.xml} and hourly CSV matrices \\
  \texttt{traffic\_generator.py} & high-level orchestration of the three pipeline stages \\
  \bottomrule
  \end{tabularx}
  \caption{Structure of the \texttt{scenario\_generator} package.}
  \label{tab:scenario-generator-concepts}
\end{table}

Table \ref{tab:scenario-generator-concepts} summarizes the main concepts of the
\texttt{scenario\_generator}. The generator takes as input : 
\begin{itemize}
  \item a SUMO road network (\texttt{.net.xml}),
  \item a Traffic Analysis Zone (TAZ) file (\texttt{taz.xml}),
  \item a TOML configuration file specifying demand and routing parameters.
\end{itemize}

It produces:
\begin{itemize}
  \item a SUMO vehicle route file (\texttt{routes.xml}) that can be used
        directly as a demand input to {\texttt{sumo}} or {\texttt{sumo-gui}},
  \item hourly OD matrices in both XML ({\texttt{out\_od.xml}}) and CSV form
        for inspection and re-use.
\end{itemize}

The pipeline consists of three sequential stages:
\begin{enumerate}
  \item \textbf{Stage 1 - Demand generation}: construct 24 hourly
        Origin\nobreakdash–Destination (OD) matrices using a gravity model
        with an optional directional commute overlay.
  \item \textbf{Stage 2 - Route library}: build a small set of candidate
        routes for each active OD pair using Yen's $k$-shortest
        paths algorithm \cite{yen_algorithm}.
  \item \textbf{Stage 3 - Vehicle instantiation}: assign a departure time
        and a concrete route to every trip.
\end{enumerate}

\subsubsection{Stage 1 : Demand generation}

\paragraph{Traffic Analysis Zones}
The spatial unit of the model is the \emph{Traffic Analysis Zone} (TAZ),
a standard construct in transport demand modelling that aggregates individual
trip origins and destinations into discrete geographic areas.  
\href{https://sumo.dlr.de/docs/Contributed/SUMOPy/Demand/Zone_To_Zone.html}{TAZs in SUMO} 
are defined as polygons that can be associated with the road network.
Each TAZ is defined by the set of SUMO road edges it contains.

A scalar \emph{zone mass} $m_i$ is assigned to each TAZ as the capacity of its drivable
edge lengths : 
\begin{equation}
  m_i = \sum_{e \in \mathcal{E}_i} \ell(e) \cdot n(e)
  \label{eq:zone_mass}
\end{equation}
where $\ell(e)$ is the length of edge $e$ in meters and $n(e)$ is its lane count. 

Equation \ref{eq:zone_mass} is a \emph{deliberate approximation of the network structure}
adopted in the absence of socio-economic calibration data. In the classical four-step
transport model, trips productions and attractions are estimated using socio-economic data
\cite{modelling_transport} but none of those are available in the present synthetic
setting. Lane kilometers \footnote{A lane kilometer is a unit of length that is equal to 
one kilometer of road with one lane.} are used here as a surrogate because they are a 
directly measurable network attribute and because major road corridors tend to account for 
a disproportionate share of traffic volume relative to their share of network length.
The resulting $m_i$ has therefore to be replaced by observed productions and attractions
whenever such data become available. 

\paragraph{Temporal Demand Profile}
\label{par:temporal_demand_profile}

The total daily demand $N$ is distributed across the 24 hours of the
simulation using a \emph{two\nobreakdash-Gaussian profile}:
\begin{equation}
  p_h = \frac{a_m\,g(h;\,\mu_m,\sigma) + a_e\,g(h;\,\mu_e,\sigma)}
             {\displaystyle\sum_{h'=0}^{23}
              \bigl[a_m\,g(h';\,\mu_m,\sigma)
                  + a_e\,g(h';\,\mu_e,\sigma)\bigr]},
  \label{eq:twopeak_gaussian}
\end{equation}
where 
\begin{equation}
  g(h;\,\mu,\sigma) = \exp\!\left(-\frac{(h-\mu)^2}{2\sigma^2}\right),
\end{equation}
$\mu_m$ and $\mu_e$ are the morning and evening peak hours respectively, and
$\sigma$ controls the temporal spread.  The amplitudes $a_m$ and $a_e$ allow
the two peaks to differ in magnitude, the default configuration sets
$a_e > a_m$ to reflect the slightly larger afternoon peak typically observed
in urban network \cite{gaussian_book}.

The continuous profile $p_h$ is converted into hourly trip counts $n_h$ using the
largest-remainder method which guarantees $\sum_h n_h = N$ while preserving profile
proportions as closely as possible under the integer constraint.

\begin{figure}[htbp]
  \centering
  \includegraphics[width=0.92\textwidth]{plot_hourly_demand.pdf}
  \caption{Hourly trip distribution produced by the two-Gaussian profile
           (Equation~\ref{eq:twopeak_gaussian}) for the test scenario.
           The orange bar marks the evening peak, which is slightly larger
           than the morning peak ($a_e > a_m$).}
  \label{fig:hourly_demand}
\end{figure}

\paragraph{Gravity OD model}
\label{par:gravity_od_model}

The spatial distribution of trips across TAZ pairs is governed by a \emph{singly-constrained
gravity model} \cite{gravity_model}. The \textbf{unnormalized} weight of the OD pair $(i,j)$ is :
\begin{equation}
  w_{ij} = m_i \cdot m_j \cdot \exp(-\beta\, c_{ij}), \quad i \neq j,
  \label{eq:gravity_model}
\end{equation}
where $c_{ij}$ is the travel-cost impedance between zones $i$ and $j$ (shortest path length
in meters) and $\beta \geq 0$ is the distance decay parameter\footnote{A larger $\beta$
means trips decay faster with distance/cost (people are more reluctant to travel far), a smaller $\beta$
means the distribution is more spatially spread out. Wilson discusses this in the context of 
calibrating the model to observed data.\cite{gravity_model}}.

Pairs with infinite or negative impedance are assigned to zero weight. The \textbf{normalized}
\emph{OD share matrix} is obtained as : 
\begin{equation}
  S_{ij} \;=\; \underbrace{\frac{m_i}{\sum_{k} m_k}}_{\text{origin share }O_i}
              \;\cdot\;
              \underbrace{\frac{m_j\,\exp(-\beta c_{ij})}
                               {\sum_{\ell} m_\ell\,\exp(-\beta c_{i\ell})}}_{\text{conditional destination probability }p(j\mid i)},
  \label{eq:od_share}
\end{equation}
so that, by construction, $\sum_j S_{ij} = m_i / \sum_k m_k$ and
$\sum_{ij} S_{ij} = 1$.

For each hour $h$, the share matrix is scaled to the corresponding integer total $n_h$ 
via the largest-remainder method, yielding an exact integer OD matrix $T^h$ with
$\sum_{ij} T^h_{ij} = n_h$.

The model is \emph{singly-constrained on the production side}. Although the
weight $w_{ij} = m_i m_j \exp(-\beta c_{ij})$ depends symmetrically on origin
and destination masses, the normalisation in Equation~\ref{eq:od_share} enforces
the row-sum constraint
\begin{equation*}
  \sum_j S_{ij} \;=\; \frac{m_i}{\sum_k m_k}
\end{equation*}
while leaving the column sums $\sum_i S_{ij}$ unconstrained. In Wilson's
balancing-factor formulation \cite{ortuzar_modelling_transport}, this
corresponds to enforcing only the origin-side factor $A_i$ and fixing the
destination-side factor $B_j \equiv 1$. A doubly-constrained model would
additionally enforce observed zonal attractions and would require both
balancing factors to be recovered by iterative proportional fitting; zonal
attractions are not observed in the present synthetic setting, so enforcing
them would introduce unidentified parameters. The singly-constrained
formulation is therefore the parsimonious choice that uses only the
information actually contained in the network (zone masses and impedances)
and leaves destination totals free to respond to the network structure
through $\exp(-\beta c_{ij})$.

\begin{figure}[htbp]
  \centering
  \includegraphics[width=\textwidth]{plot_od_heatmap.pdf}
  \caption{Daily OD demand matrix $\sum_h T^h$ for the test scenario (20 TAZs,
           10\,000 daily trips). Rows are origin zones, columns are destination
           zones. The cell colour and annotation encode the total number of trips.
           The diagonal is zero by construction (no intra demand). Dense
           clusters correspond to adjacent zones with high lane-kilometer mass.}
  \label{fig:od_heatmap}
\end{figure}

\paragraph{Commute Overlay (Optional)}
\label{par:commute_overlay}

When residential weights $r_i$ and employment weights $e_i$ are provided
for each zone, a directional commute bias is introduced into the hourly
OD shares. Two commute matrices are constructed following the same
gravity structure as Equation~\ref{eq:gravity_model}: one for the morning
peak, capturing residential-to-employment flows, and one for the evening
peak, capturing the reverse direction:
\begin{align}
  w^{\mathrm{AM}}_{ij} &= m_i\, m_j\, f(c_{ij})\, r_i\, e_j, \\
  w^{\mathrm{PM}}_{ij} &= m_i\, m_j\, f(c_{ij})\, e_i\, r_j,
\end{align}
where $f(c_{ij}) = \exp(-\beta_c\, c_{ij})$ is a distance-decay function
that discounts longer or costlier trips. In the AM matrix, the additional
factors $r_i$ and $e_j$ amplify flows leaving residential zones and
arriving at employment zones, the PM matrix reverses this logic by
weighting departures from employment zones ($e_i$) toward residential
destinations ($r_j$).

These directional matrices are then blended with the base OD share matrix
through a time-varying mixture. For each hour $h$, the blended share is:
\begin{equation}
  S^h_{ij} \;\propto\; S^{\mathrm{base}}_{ij}
  + \lambda\, g(h;\,\mu_m,\,\sigma)\, S^{\mathrm{AM}}_{ij}
  + \lambda\, g(h;\,\mu_e,\,\sigma)\, S^{\mathrm{PM}}_{ij},
  \label{eq:commute_blend}
\end{equation}
where $g(h;\,\mu,\,\sigma)$ is a Gaussian kernel centred on hour $\mu$
with spread $\sigma$, and $\lambda$ (\texttt{peak\_weight}) governs the
overall strength of the commute signal relative to the background. The AM
component peaks around the morning rush hour $\mu_m$ and fades outside
it, the PM component does the same around the evening rush hour $\mu_e$.
Setting $\lambda = 0$ reduces the blended share to the base model alone.
The resulting shares are normalised before being applied to the trip
scaling step.

\begin{figure}[htbp]
  \centering
  \includegraphics[width=\textwidth]{peak_directionality.png}
  \caption{Effect of the commute overlay on the hourly OD matrices.
           Left column: base gravity model ($\lambda = 0$). Centre column:
           model with commute overlay ($\lambda > 0$). Right column:
           difference (centre $-$ left). Top row: morning peak hour
           ($h = 8$). Bottom row: evening peak hour ($h = 17$). Red cells
           in the difference maps indicate flows amplified by the overlay,
           bylue cells indicate flows attenuated. The AM overlay concentrates
           demand on residential$\to$employment corridors, the PM overlay
           reverses this pattern.}
  \label{fig:commute_overlay}
\end{figure}

\subsubsection{Stage 2 : Route library}
\label{sssec:route_library}

\paragraph{Edge expanded network}
\label{par:edge_expanded_network}

Routing is performed on an \emph{edge-expanded} directed graph in which SUMO
road edges serve as nodes. A directed arc connects edge $u$ to edge $v$
whenever $v$ is reachable from $u$ through a SUMO connection, with arc weight
equal to the length of $v$, path cost is therefore the sum of edge lengths in
metres.

The edge-expanded representation is preferred to a junction-level (node-based)
graph for two reasons specific to this pipeline. First, SUMO's demand format
(\texttt{routes.xml}) requires routes to be expressed as ordered edge
sequences, an edge-level graph produces paths that can be serialised directly,
whereas a node-level graph would require an additional junction-to-edge
conversion. Second, the edge-expanded graph natively encodes turning
feasibility, an arc $(u, v)$ exists if and only if SUMO permits the
corresponding connection, so every path returned by the router is
guaranteed to be physically drivable without any posterior feasibility check.
Internal junction edges \footnote{SUMO ids beginning with \texttt{:}} and edges that
do not permit private-vehicle access are excluded at graph construction time.


\paragraph{Impedance Estimation}
\label{par:impedance_estimation}

The gravity model in Stage 1 requires a TAZ-to-TAZ impedance matrix. Computing exact 
center-to-center shortest paths would require all-pairs Dijkstra on a large graph,
which is unnecessarily expensive for the present purpose.  Instead, for each
directed zone pair $(i,j)$, a small number of (origin-edge, destination-edge) combinations 
is sampled from the two zones and the minimum observed single-source shortest-path cost
is taken as the representative impedance.  The approximation is decreasing in 
geographic proximity, which is sufficient for the gravity model to produce a spatially 
plausible OD distribution.

\paragraph{$k$-Shortest Candidate Routes}
\label{par:k_shortest_candidate_routes}

For each active OD pair (i.e.\ $T^h_{ij} > 0$ for at least one hour $h$),
a set of at most $k$ candidate routes is computed as follows:
\begin{enumerate}
  \item Sample $s$ concrete (origin-edge, destination-edge)
        pairs, where edges are drawn proportionally to zone mass
        (Equation~\eqref{eq:zone_mass}).
  \item For each pair, enumerate up to $k$ shortest simple paths using
        Yen's algorithm \cite{yen_algorithm}
  \item Deduplicate across all samples by edge sequence, if the same
        sequence appears more than once, retain the lowest-cost
        instance.
  \item Return at most $k$ candidates sorted by cost.
\end{enumerate}

\subsubsection{Stage 3 : Vehicle Instantiation}
\label{sssec:vehicle_instantiation}

Once the hourly OD matrices and the candidate route library are available,
every trip in the demand model is turned into a concrete SUMO vehicle
definition.  Each trip requires two random draws: a departure time and a
route.

\paragraph{Departure Time}
\label{par:departure_time}

Trips within each one-hour interval $[b, b+3600)$ are assumed
to depart uniformly over that interval.  For an OD cell containing $n$ trips,
departure times are drawn independently as:
\begin{equation}
  t_1,\ldots,t_n \;\overset{\mathrm{i.i.d.}}{\sim}\; \mathcal{U}(b,\; b+3600).
\end{equation}
This is the natural maximum-entropy assumption when only hourly
totals are known: given no information about the exact timing of individual
trips within an hour, all departure instants in $[b, b+3600)$ are treated as
equally likely.

\paragraph{Route Selection}
\label{par:route_selection}

Each trip is assigned one route selected \emph{uniformly at random} from the
$k$ precomputed candidates for its OD pair:
\begin{equation}
  P(r \mid \mathcal{R}) = \frac{1}{|\mathcal{R}|}, \quad r \in \mathcal{R}.
  \label{eq:uniform_choice}
\end{equation}
This is again the maximum-entropy choice in the absence of
driver-preference data: if no reason is given to favour one
candidate over another, all candidates are equally likely.  The approach
ensures that, on average, the generated traffic load is spread across the
full set of precomputed alternatives rather than being
concentrated on a single route per OD pair.

The two draws (departure time and route) are performed with a single shared
random-number generator whose seed is fixed in the configuration
file, making the full scenario exactly reproducible.

\newpage
\subsubsection{Summary of Parameters}
\label{sssec:sg_params}

\begin{table}[ht]
  \centering
  \caption{Key parameters of the scenario generator and their roles.}
  \label{tab:sg_params}
  \begin{tabular}{lll}
    \toprule
    Parameter & Stage & Effect \\
    \midrule
    {\ttfamily daily\_total}                & 1 & Total daily trip count $N$ \\
    {\ttfamily distance\_decay\_beta} $\beta$ & 1 & Gravity model decay; higher $\Rightarrow$ shorter trips \\
    {\ttfamily morning\_peak\_hour} $\mu_m$ & 1 & AM peak centre (hours) \\
    {\ttfamily evening\_peak\_hour} $\mu_e$ & 1 & PM peak centre (hours) \\
    {\ttfamily sigma\_hours} $\sigma$       & 1 & Temporal spread of peaks \\
    {\ttfamily peak\_weight} $\lambda$      & 1 & Commute overlay strength \\
    {\ttfamily k\_paths} $k$               & 2   & Number of candidate routes per OD pair \\
    {\ttfamily endpoint\_samples\_per\_od} $s$ & 1\,\&\,2 & Accuracy vs.\ runtime trade-off \\
    \bottomrule
  \end{tabular}
\end{table}